% =====================================================================
% Runbook Reproduksi & Operasional SmartCityApps
% Standalone document — compile with: latexmk -pdf runbook.tex
% =====================================================================
\documentclass[11pt,a4paper]{article}

\usepackage[T1]{fontenc}
\usepackage[utf8]{inputenc}
\usepackage[a4paper,margin=2.4cm]{geometry}
\usepackage{lmodern}
\usepackage{xcolor}
\usepackage{listings}
\usepackage{booktabs}
\usepackage{enumitem}
\usepackage{titlesec}
\usepackage{fancyhdr}
\usepackage[hidelinks,colorlinks=true,linkcolor=black!75,urlcolor=blue!55!black]{hyperref}

% ---------- colours ----------
\definecolor{accent}{HTML}{2E4053}
\definecolor{codebg}{HTML}{F4F6F7}
\definecolor{codeframe}{HTML}{D5DBDB}
\definecolor{kw}{HTML}{1E5AA8}
\definecolor{cmt}{HTML}{7B7D7D}
\definecolor{stagebg}{HTML}{EAF2F8}

% ---------- headings ----------
\titleformat{\section}{\Large\bfseries\color{accent}}{\thesection}{0.6em}{}
\titleformat{\subsection}{\large\bfseries\color{accent!90}}{\thesubsection}{0.6em}{}
\titlespacing*{\section}{0pt}{1.4em}{0.6em}

% ---------- header / footer ----------
\pagestyle{fancy}\fancyhf{}
\setlength{\headheight}{14pt}
\lhead{\small\color{accent}\textbf{Runbook SmartCityApps}}
\rhead{\small\color{accent}\thepage}
\renewcommand{\headrulewidth}{0.4pt}

% ---------- code style ----------
\lstdefinestyle{sh}{
  language=bash,
  basicstyle=\ttfamily\small,
  backgroundcolor=\color{codebg},
  frame=single, rulecolor=\color{codeframe},
  framesep=6pt, xleftmargin=8pt, xrightmargin=4pt,
  breaklines=true, breakatwhitespace=false,
  postbreak=\mbox{\textcolor{cmt}{$\hookrightarrow$}\space},
  showstringspaces=false,
  keywordstyle=\color{kw}\bfseries,
  commentstyle=\color{cmt}\itshape,
  morecomment=[l]{\#},
  morekeywords={python,python3,pip,uvicorn,curl,docker,ngrok,flutter,adb,cd,source,jupyter,latexmk,bibtex},
  aboveskip=8pt, belowskip=8pt,
}
\lstset{style=sh}

% command shown inline
\newcommand{\code}[1]{\texttt{\small #1}}

% stage banner
\newcommand{\stage}[2]{%
  \par\medskip\noindent\colorbox{stagebg}{\parbox{\dimexpr\linewidth-2\fboxsep}{%
  \textbf{\color{accent}#1}\;\;#2}}\par\smallskip}

\begin{document}

% =====================================================================
\begin{titlepage}
\centering
\vspace*{2.2cm}
{\Huge\bfseries\color{accent} Runbook Reproduksi \&\\[0.25em] Operasional SmartCityApps\par}
\vspace{0.8cm}
{\large Dari Proses Analisis Data hingga Inferensi Server dan Aplikasi Mobile\par}
\vspace{0.4cm}
{\large\itshape Pipeline Klasifikasi Multimodal CRM Jakarta\\(DINOv3-large + multilingual-E5-large + CatBoost+PGS)\par}
\vfill
\begin{tabular}{rl}
\textbf{Repositori} & \code{smartCityReport/}\\
\textbf{Server inferensi} & FastAPI \code{serve\_model.py} (\code{/health}, \code{/predict})\\
\textbf{Klien} & Flutter Android \code{smart\_city\_reporter\_app/}\\
\textbf{Backend data} & Supabase (PostgreSQL, Storage, Auth)\\
\textbf{Publikasi} & ngrok tunnel\\
\end{tabular}
\vspace{1.2cm}

{\small Dokumen ini bersifat operasional (\emph{runbook}); jalankan tahap secara berurutan.}
\vfill
\end{titlepage}

\tableofcontents
\newpage

% =====================================================================
\section{Ikhtisar Alur}
% =====================================================================
Runbook ini menjalankan keseluruhan sistem secara berurutan, mulai dari analisis
data mentah sampai aplikasi mobile mengonsumsi hasil prediksi. Sembilan tahap
utama dirangkum pada Tabel~\ref{tab:overview}.

\begin{table}[h]
\centering
\renewcommand{\arraystretch}{1.25}
\begin{tabular}{@{}clp{8.4cm}@{}}
\toprule
\textbf{Tahap} & \textbf{Nama} & \textbf{Keluaran} \\
\midrule
A & Analisis \& Persiapan Data & \code{train/val/test.csv}, statistik kelas \\
B & Ekstraksi \emph{Embedding} & \emph{cache} fitur citra \& teks (1024-d) \\
C & Pelatihan Fusi CatBoost+PGS & \code{.cbm} model + metrik validasi \\
D & Ekspor ONNX \& Paritas & artefak \code{artifacts/export/} \\
E & Server Inferensi FastAPI & layanan \code{:8000} (\code{/predict}) \\
F & Kontainerisasi Docker & image \code{crm-multimodal:latest} \\
G & Publikasi ngrok & URL HTTPS publik \\
H & Aplikasi Flutter & APK / sesi \code{flutter run} \\
I & Verifikasi End-to-End & laporan tersimpan di Supabase \\
\bottomrule
\end{tabular}
\caption{Sembilan tahap runbook dari analisis hingga konsumsi aplikasi.}
\label{tab:overview}
\end{table}

% =====================================================================
\section{Prasyarat \& Lingkungan}
% =====================================================================
\begin{itemize}[leftmargin=1.4em,itemsep=2pt]
  \item \textbf{Python} 3.11 (server \& notebook).
  \item \textbf{Flutter} SDK (channel stabil) + Android SDK / perangkat Android.
  \item \textbf{ngrok} (akun gratis cukup) untuk publikasi server.
  \item \textbf{Docker} (opsional, untuk jalur kontainer/GPU).
  \item \textbf{GPU} (opsional): pelatihan \& ekstraksi \emph{embedding} dijalankan di RunPod; inferensi dapat berjalan di CPU (mis. macOS / Apple Silicon).
\end{itemize}

\stage{Catatan platform.}{Pelatihan dan ekspor (Tahap A--D) dilakukan pada pod RunPod ber-GPU. Penyajian (Tahap E--I) dapat berjalan penuh di CPU lokal menggunakan subset dependensi \code{requirements-serve.txt}.}

\subsection{Menyiapkan virtualenv (penyajian, CPU)}
\begin{lstlisting}
cd smartCityReport
python3.11 -m venv .venv
source .venv/bin/activate
pip install -U pip
pip install -r requirements-serve.txt    # fastapi, uvicorn, onnxruntime (CPU),
                                          # transformers, catboost, pillow, numpy
\end{lstlisting}

\subsection{Menyiapkan dependensi pelatihan (GPU / RunPod)}
\begin{lstlisting}
# Pada pod RunPod (CUDA):
pip install -r requirements.txt              # torch (CUDA), timm, dsb.
pip install -r notebooks/requirements_runpod.txt
\end{lstlisting}

% =====================================================================
\section{Tahap A — Analisis \& Persiapan Data}
% =====================================================================
Dataset CRM Jakarta (citra + narasi + label SKPD) dianalisis, dibersihkan, dan
dibagi secara stratifikasi. Jalankan notebook dari dalam folder \code{notebooks/}
agar \code{sys.path} relatif (\code{../src/}) ter-resolve.

\begin{lstlisting}
cd smartCityReport/notebooks
jupyter notebook        # atau: jupyter lab
\end{lstlisting}

\begin{enumerate}[leftmargin=1.6em,itemsep=3pt]
  \item \textbf{\code{01\_dataset\_exploration.ipynb}} — distribusi 9 kelas,
        contoh citra, analisis ukuran gambar, distribusi 487 SKPD mentah.
        \emph{Output:} \code{class\_distribution.png}, \code{skpd\_distribution.png}.
  \item \textbf{\code{02\_preprocessing.ipynb}} — normalisasi unicode,
        \emph{whitespace collapsing}, pembersihan label, pemetaan 487 SKPD
        $\rightarrow$ 9 instansi target.
  \item \textbf{\code{03\_split.ipynb}} — \emph{stratified split} berdasarkan
        \code{label\_id}. \emph{Output:} \code{train.csv} (43.241),
        \code{val.csv} (9.266), \code{test.csv} (9.266).
\end{enumerate}

\stage{Validasi tahap.}{Pastikan proporsi label tiap kelas konsisten lintas partisi \code{train}/\code{val}/\code{test}, dan \emph{class weights} (inverse frequency) tersimpan untuk dipakai CatBoost.}

% =====================================================================
\section{Tahap B — Ekstraksi \emph{Embedding} dengan \emph{Encoder} Beku}
% =====================================================================
Setiap citra dan narasi dilewatkan melalui \emph{encoder} beku untuk menghasilkan
vektor 1024-dimensi; tidak ada gradien yang mengalir ke \emph{encoder}.

\begin{enumerate}[leftmargin=1.6em,itemsep=3pt]
  \item \textbf{\code{04\_extract\_image\_embeddings.ipynb}} — DINOv3-large
        ($resize \rightarrow center\ crop \rightarrow$ normalisasi ImageNet
        $\rightarrow$ NCHW). \emph{Output:} \emph{cache} \code{image\_emb.npy}.
  \item \textbf{\code{05\_extract\_text\_embeddings.ipynb}} — multilingual-E5-large
        (\emph{tokenize} $\rightarrow$ \emph{mean pooling} $\rightarrow$ L2-normalize).
        \emph{Output:} \emph{cache} \code{text\_emb.npy}.
\end{enumerate}

\stage{Tujuan caching.}{Karena \emph{encoder} beku, \emph{embedding} dihitung sekali lalu dipakai ulang oleh seluruh eksperimen pelatihan CatBoost, sehingga eksperimen menjadi cepat dan reproducible.}

% =====================================================================
\section{Tahap C — Pelatihan Kepala CatBoost + PGS}
% =====================================================================
\emph{Embedding} citra dan teks dikonkatenasi (\emph{early fusion}) menjadi vektor
2048-d, lalu CatBoost dilatih untuk meminimalkan \emph{multiclass logloss}.

\begin{enumerate}[leftmargin=1.6em,itemsep=3pt]
  \item \textbf{\code{06\_fusion\_catboost\_pgs.ipynb}} — pelatihan CatBoost +
        kalibrasi \emph{Posterior Gaussian Sampling}; pemilihan model terbaik
        berbasis \emph{macro F1} pada set validasi.
        \emph{Output:} \code{artifacts/models/best\_catboost\_dinov3\_large\_mE5\_large.cbm}.
  \item \textbf{\code{07\_error\_analysis.ipynb}} / \code{07\_standalone\_fusion\_analysis.ipynb}
        — \emph{confusion matrix}, analisis kesalahan, \emph{threshold trade-off}.
  \item \textbf{\code{08\_add\_dinov3\_ablation.ipynb}}, \code{09\_advanced\_classifiers.ipynb}
        — ablasi modalitas dan pasangan \emph{encoder} ($4\times4\times2$).
\end{enumerate}

\stage{Hasil rujukan.}{Konfigurasi terbaik mencapai \emph{balanced accuracy} 0{,}8074 dan \emph{macro F1} 0{,}7747 pada 9.266 sampel uji.}

% =====================================================================
\section{Tahap D — Ekspor ONNX \& Validasi Paritas}
% =====================================================================
\begin{enumerate}[leftmargin=1.6em,itemsep=3pt]
  \item \textbf{\code{08\_export\_onnx.ipynb}} — ekspor \emph{image encoder},
        \emph{text encoder}, dan \emph{classifier head} ke ONNX.
\end{enumerate}

Artefak final yang harus tersedia sebelum penyajian:
\begin{lstlisting}
artifacts/export/
  image_encoder_dinov3_large.onnx            # 1.15 GB
  image_encoder_dinov3_large/preprocessor_config.json
  text_encoder_mE5_large/model.onnx (+ model.onnx_data ~6.4 GB)
  tokenizer/
  classifier_dinov3_large_mE5_large.onnx     # CatBoost-as-ONNX
\end{lstlisting}

\stage{Validasi paritas.}{Bandingkan probabilitas keluaran PyTorch vs ONNX pada sampel uji; selisih absolut maksimum harus di bawah ambang $10^{-5}$ agar prediksi top-1 identik.}

% =====================================================================
\section{Tahap E — Menjalankan Server Inferensi FastAPI}
% =====================================================================
Server memuat ketiga artefak ONNX sekali saat startup (\(\approx\) 7{,}6 GB di RAM),
lalu menyediakan \code{/health} dan \code{/predict}.

\begin{lstlisting}
cd smartCityReport
source .venv/bin/activate
uvicorn serve_model:app --host 0.0.0.0 --port 8000 --workers 1
\end{lstlisting}

\stage{Penting.}{Gunakan \code{-{}-workers 1}: setiap worker memuat penuh \(\approx\)7{,}6 GB. \emph{Cold start} lambat (memuat \emph{text encoder} 6{,}4 GB) — lakukan satu \code{/health} dan satu \code{/predict} pemanasan sebelum demo.}

\subsection{Uji \emph{health}}
\begin{lstlisting}
curl -s http://127.0.0.1:8000/health | python3 -m json.tool
# -> {"status":"ok", "image_encoder":"dinov3_large",
#     "text_encoder":"mE5_large", "classes":[...9...], "gpu_available":false}
\end{lstlisting}

\subsection{Uji \emph{predict} (multipart)}
\begin{lstlisting}
curl -s -X POST http://127.0.0.1:8000/predict \
  -F "image=@/path/to/photo.jpg" \
  -F "laporan=ada lubang besar di Jalan Sudirman" \
  -F "latitude=-6.2088" -F "longitude=106.8456" | python3 -m json.tool
# -> predicted_dinas, confidence, all_probabilities,
#    review_required, assignment{agency_name, distance_meters, ...}
\end{lstlisting}

% =====================================================================
\section{Tahap F — Kontainerisasi Docker (opsional, GPU)}
% =====================================================================
Jalur ini memakai \code{Dockerfile} berbasis CUDA 12.1 + \code{onnxruntime-gpu};
membutuhkan host ber-GPU NVIDIA (mis. pod RunPod), bukan macOS.

\begin{lstlisting}
cd smartCityReport
docker compose up --build          # build image + mount artifacts/export (read-only)
# image: crm-multimodal:latest, port 8000, GPU direservasi via compose.yaml
\end{lstlisting}

\stage{Catatan artefak.}{\code{compose.yaml} me-mount \code{artifacts/export} sebagai volume \emph{read-only} sehingga artefak 7{,}6 GB tidak perlu di-\emph{bake} ke dalam image.}

% =====================================================================
\section{Tahap G — Publikasi Server via ngrok}
% =====================================================================
Agar perangkat Android dapat menjangkau server melalui internet:

\begin{lstlisting}
brew install ngrok                       # atau unduh dari ngrok.com
ngrok config add-authtoken <TOKEN>       # sekali saja (akun gratis)
ngrok http 8000
# -> Forwarding  https://abc123.ngrok-free.app -> http://localhost:8000
\end{lstlisting}

Verifikasi dari luar jaringan:
\begin{lstlisting}
curl -s https://abc123.ngrok-free.app/health
\end{lstlisting}

\stage{Peringatan.}{URL ngrok gratis berubah setiap restart. Catat URL terbaru dan teruskan ke aplikasi (Tahap H) sebelum demo.}

% =====================================================================
\section{Tahap H — Menjalankan Aplikasi Flutter}
% =====================================================================
Aplikasi membaca alamat server dari \emph{dart-define} \code{CRM\_API\_URL}
(default \code{http://10.0.2.2:8000} untuk emulator Android).

\begin{lstlisting}
cd smartCityReport/smart_city_reporter_app
flutter pub get
flutter run \
  --dart-define=CRM_API_URL=https://abc123.ngrok-free.app \
  --dart-define=SUPABASE_URL=<supabase_url> \
  --dart-define=SUPABASE_ANON_KEY=<anon_key>
\end{lstlisting}

\subsection{Alternatif tanpa ngrok}
\begin{itemize}[leftmargin=1.4em,itemsep=2pt]
  \item \textbf{Emulator:} biarkan default \code{http://10.0.2.2:8000} (alias localhost host).
  \item \textbf{HP via USB:} \code{adb reverse tcp:8000 tcp:8000}, lalu
        \code{CRM\_API\_URL=http://localhost:8000}.
\end{itemize}

% =====================================================================
\section{Tahap I — Verifikasi End-to-End}
% =====================================================================
\begin{enumerate}[leftmargin=1.6em,itemsep=2pt]
  \item Di aplikasi: ambil foto $\rightarrow$ tulis narasi $\rightarrow$ izinkan lokasi $\rightarrow$ kirim.
  \item Layar tinjauan AI menampilkan \textbf{instansi prediksi, confidence, top-3, instansi terdekat} (hasil nyata dari server).
  \item Koreksi kategori bila perlu $\rightarrow$ simpan.
  \item Laporan tersimpan di Supabase (tabel \code{reports}) + foto di \emph{bucket} \code{report-images}.
  \item Operator Instansi memverifikasi; status \code{resolved} mewajibkan foto bukti penyelesaian.
\end{enumerate}

% =====================================================================
\section{Troubleshooting}
% =====================================================================
\begin{table}[h]
\centering
\renewcommand{\arraystretch}{1.3}
\begin{tabular}{@{}p{5.2cm}p{9.0cm}@{}}
\toprule
\textbf{Gejala} & \textbf{Penyebab / Solusi} \\
\midrule
Server OOM / lambat saat start & RAM kurang untuk 7{,}6 GB; gunakan \code{-{}-workers 1}, tutup aplikasi lain. \\
\code{Invalid image} pada \code{/predict} & Berkas bukan gambar valid; kirim JPG/PNG via \code{-F "image=@..."}. \\
Aplikasi tak terhubung & URL ngrok kedaluwarsa; perbarui \code{CRM\_API\_URL} dan \code{flutter run} ulang. \\
Emulator gagal \code{localhost} & Gunakan \code{10.0.2.2} (emulator) atau \code{adb reverse} (HP fisik). \\
Paritas ONNX gagal ($>10^{-5}$) & Periksa opset / operator ekspor pada \code{08\_export\_onnx.ipynb}. \\
Email verifikasi tidak masuk & SMTP bawaan Supabase ber-\emph{rate limit}; gunakan SMTP khusus untuk volume lebih besar. \\
\bottomrule
\end{tabular}
\caption{Masalah umum dan penanganannya.}
\end{table}

% =====================================================================
\section{Lampiran — Peta Berkas Penting}
% =====================================================================
\begin{lstlisting}
smartCityReport/
  serve_model.py               # FastAPI: /health, /predict, routing Haversine
  Dockerfile, compose.yaml     # jalur kontainer GPU
  requirements-serve.txt       # subset penyajian (CPU)
  requirements.txt             # pelatihan (CUDA)
  notebooks/01..09_*.ipynb     # analisis -> pelatihan -> ekspor
  src/crm/encoders/{image,text}.py  # pembungkus encoder ONNX
  src/crm/fusion.py, pgs.py    # early fusion + kalibrasi PGS
  artifacts/export/            # artefak ONNX yang dilayani
  smart_city_reporter_app/     # aplikasi Flutter (klien)
    lib/core/config/app_config.dart   # CRM_API_URL, SUPABASE_URL, ANON_KEY
\end{lstlisting}

\vfill
\begin{center}\small\color{cmt}
Runbook ini melengkapi Bab 3 (Metode Penelitian) dan Bab 4 (Hasil) skripsi.
\end{center}

\end{document}
